\section{Producing and symmetrizing the marked campaign}
\label{sec:marked-campaign-producer}

This section supplies the campaign used by the fixed-stencil construction.
It has two parts. First, the equation ideal is read by microscopic link
observables whose Taylor errors lie strictly above the requested grade.
Second, finite-dimensional linear algebra combines these reads into one
campaign, and finite group averages impose the signed and real symmetries.
The filtration throughout counts factors that are color coordinates of
based curvature slots; it is not a count of raw link entries.

\subsection{The microscopic filtration and its Taylor data}

Fix a gauge group, a holomorphic link action, and a cutoff degree $d$. A
slot $S$ is one of the based, transported plaquette-curvature readouts on
the cover lattice. For a real basis vector $e_a$ of the color space, write
\[
              x_{S,a}(U)= e_a\mathbin{\cdot} S(U).
\]
Let $\mathcal P_{\geq r}$ be the complex span of finite products of at
least $r$ such coordinates, with the empty product allowed when $r=0$.
Thus a single slot coordinate has filtration degree one even though, as a
function of link matrices, it is a nonlinear polynomial or analytic
expression.

For a label $p$ of grade $m=(r,s)$, a marked read consists of a field
function $f$ and a polynomial $q$ with three properties. The function is
gauge invariant, has at least $s$ slot factors in the span sense above,
and its link-action derivatives through order $d$ agree with those of $q$.
The correction $q-p$ has strictly larger grade stamp
\[
                         \operatorname{st}(r,s)=(s,r)
\]
in lexicographic order. When $p$ lies in the Maxwell equation ideal, the
read $f$ is additionally microscopic: its moment against every sufficiently
large-volume spectator which is local away from a fixed anchor vanishes.
These are precisely the properties needed for later stencil pricing and
for the exact equation-ideal payment.

The elementary analytic input is a first-order germ for each slot. Put
$\operatorname{Snap}_d$ for the finite-dimensional space of radial link
jets through degree $d$. For every slot $S$, the function
\[
          J\longmapsto S\bigl(U_d(J)\bigr),\qquad J\in\operatorname{Snap}_d,
\]
is real analytic near zero, transforms by the adjoint action under a
constant gauge transformation at the base point, vanishes at zero, and
has derivative equal to the linearized plaquette curvature. To see the
analyticity directly, on one link the connection coefficient is a linear
function $A_J(t)$ of $J$. On the Banach space $C([0,1],\mathfrak A)$ of
continuous paths in the finite-dimensional matrix algebra $\mathfrak A$
containing the group representation, the Volterra operator
\[
          (T_Ju)(t)=\int_0^t u(s)A_J(s)\,ds
\]
has norm at most $\sup_s\|A_J(s)\|$. In a neighborhood of $J=0$ this is
less than one, so the unique solution of the link equation is
\[
                 u_J=(I-T_J)^{-1}{\bf 1}
                     =\sum_{n\geq0}T_J^n{\bf 1}.
\]
The series converges locally uniformly in the finite-dimensional parameter
$J$, and hence is analytic there. Uniqueness identifies it with the given
link transport. A slot is obtained by a finite product of these transports
and their inverses, followed by the plaquette curvature map; it is therefore
analytic as well. Constant-gauge covariance follows by conjugating the ODE,
and differentiating at $J=0$ gives the signed sum of the edge tangents around
the plaquette, which is the linearized curvature. This proves the required
slot germ without imposing any regularity on a configuration space of
infinite link fields.

\subsection{Avoiding one edge while spanning Taylor tests}

The next finite-dimensional fact is the source of the locality needed for
products. Regard $\operatorname{Snap}_d$ itself as the real vector space
$\mathcal T_d$ of radial jets. For a slot $S$, let
\[
       B_S=D_0\bigl(J\mapsto S(U_d(J))\bigr):
                  \operatorname{Snap}_d\to\mathfrak g.
\]
If $e$ is a prescribed cover edge, call a slot $S$ safe for $e$ if its
spatial support avoids $e$ along the translated path used to define $S$.

\begin{lemma}[Safe Taylor tests span]\label{lem:marked-safe-span}
For each edge $e$, the real linear span of the component functionals of
$B_S$, over slots $S$ safe for $e$, is all of $\mathcal T_d^*$.
\end{lemma}
\begin{proof}
For $d<2$ the radial-jet space has no Taylor components and the claim is
immediate. Assume $d\geq2$.
Take a jet $A\in\operatorname{Snap}_d$ annihilated by every listed
functional. For every pair of axes $i<j$, the safe-slot geometry excludes
only a finite set of plaquette base points. At each other base point it
supplies a detour transport path avoiding $e$. The hypothesis says the
resulting slot derivative on $A$ is zero. The slot-germ identity therefore
gives the full sampled curl equation below. Here
$C_{ij}(A_\ell)(x,t)$ denotes the axes-$(i,j)$ finite-mesh curl component
of the degree-$\ell$ Taylor field at position $x$ and mesh $t$; it is
homogeneous under simultaneous scaling of $x$ and $t$.
\[
       \sum_{\ell=0}^{d-2}C_{ij}(A_\ell)(x_y,1)=0
\]
at every lattice point $y$ outside the finite exceptional set. We now
induct downwards: once shells above $k$ have been shown to vanish, this
same equation reduces to its sum over $0\leq\ell\leq k$. For $t>0$ define
\[
 F_k(x,t)=\sum_{\ell=0}^{k}
              t^{k-\ell}C_{ij}(A_\ell)(x,t).
\]
Homogeneity of the degree-$\ell$ curl component gives
\[
 F_k(t x_y,t)=t^k\sum_{\ell=0}^k C_{ij}(A_\ell)(x_y,1)=0.
\]
Choose lattice sites $y_t$ with $t x_{y_t}\to x$ as $t\downarrow0$.
For $x\ne0$ the
finite exceptional set is avoided for all sufficiently small $t$.
Continuity of $F_k$ gives
$F_k(x,0)=C_{ij}(A_k)(x,0)=0$, since every lower-shell
term has a positive power of $t$ at $t=0$. Continuity extends this
identity to $x=0$. Thus the top coefficient has zero curl. For a
homogeneous radial-gauge potential of degree $k+1$, this forces the
potential to vanish directly. Differentiate its radial-gauge identity
$\sum_jx_jA_j(x)=0$ in coordinate $i$ and use curl-freeness and Euler's
identity:
\[
 0=A_i+\sum_jx_j\partial_iA_j
   =A_i+\sum_jx_j\partial_jA_i=(k+2)A_i.
\]
Hence $A_i=0$ for every coordinate, so the top shell vanishes. Remove it
from the sampled equation and repeat with the next lower index. Descending
induction proves $A=0$. In a finite-dimensional space a family of
functionals spans the dual exactly when its common kernel is zero, which
proves the claim.
\end{proof}

The slot linearization in this proof is concrete. For a based plaquette
readout it is the oriented sum of the four link tangents, transported to
the base point; the ODE calculation above proves this formula for the
actual link action. Only finitely many plaquette base points are excluded;
at each remaining point the detour path is chosen separately. The radial
argument therefore applies to every Taylor test, not just
to tests supported on a generic set.

\begin{lemma}[Equation-factor and spectator reads]\label{lem:marked-factor-read}
Let $a$ be a homogeneous polynomial in the equation ideal, of grade
$m=(r,s)$ with $r\leq d$. Then $a$ is a finite linear combination of
terms
\[
                    \ell\,\prod_{j=1}^k z_j,
\]
where $\ell$ is a homogeneous linear label functional annihilating the
Maxwell solution-potential subspace and each $z_j$ is a homogeneous linear
label atom. Their field reads are assembled from slot-coordinate
observables. Each term has a raw read $f$ and polynomial $q$ with
$q-a$ of strictly larger grade stamp, $f\in\mathcal P_{\geq s}$, and
$f$ microscopic. Individual terms need not be gauge invariant.
\end{lemma}
\begin{proof}
Write $\mathcal E_j$ for a projection onto the space of radial potentials
obtained from on-shell Maxwell curvatures at order $j$, and let $\mathcal E$
act coefficientwise. The equation
ideal is the kernel of $p\mapsto \operatorname{changed}(\mathcal E)p$.
If $p$ belongs to that ideal, then
\[
                 p=p-\operatorname{changed}(\mathcal E)p.
\]
Expand $p$ into monomials before applying the substitution. For a literal
slot atom $u$, split $u=(u-\mathcal E u)+\mathcal E u$. The first
summand annihilates every Maxwell solution potential. For a product
$u_1\cdots u_k$, use
the telescoping identity
\[
 \prod_{j=1}^k u_j-\prod_{j=1}^k\mathcal E u_j
 =\sum_{i=1}^k (u_i-\mathcal E u_i)
       \Bigl(\prod_{j<i}\mathcal E u_j\Bigr)
       \Bigl(\prod_{j>i}u_j\Bigr).
\]
The projection preserves the radial grade. Thus every summand has one
solution-annihilating linear factor and a finite spectator product, with the
original total grade. This proves the displayed factorization directly
from the equation-ideal definition; no invariant-factor assumption is
made.

Fix one generator $\ell\prod_jz_j$. The functional $\ell$ annihilates the Maxwell solution-potential subspace.
The required duality statement concerns only the top Taylor layer. For a
homogeneous potential $A_\ell$, define its transverse plaquette difference by
\[
 T_\ell(p,b;x)=N_\ell(p,b;x,1)-N_\ell(p,b;x-e_b,1),
\]
where $N_\ell$ is the sampled curl on the plaquette based at $p$. Apply the
plaquette incidence sign in direction $\nu$ to $T_\ell$ to obtain the
transverse force density. The corresponding line term
$L_\ell(\nu,p;x,t)$ is the same incidence sum of the oriented arc; each arc
is the difference of the two ordered triple averages of
$\partial_b(\partial_a A_c-\partial_c A_a)$ along the plaquette edges.
The source's scaling and zero-mesh identities are
\[
 t L_\ell(\nu,p;t x,t)=t^\ell T_\ell(\nu,p;x),\qquad
 \sum_p L_\ell(\nu,p;x,0)
       =-\sum_b\partial_b\operatorname{curl}(A_\ell)_{b\nu}(x).
\]
If every positive path-force test vanishes through order $j+2$, then for
every permitted lattice site the finite force sum vanishes. Its line form is
\[
 F_j(x,t)=\sum_{\ell=0}^j t^{j-\ell}
                    \sum_p L_\ell(\nu,p;x,t).
\]
The scaling identity shows that $tF_j(t x_y,t)$ is $t^j$ times the
vanishing finite force sum. Since $t>0$, this implies
$F_j(t x_y,t)=0$. For $x\ne0$, choose lattice sites $y_t$ with
$t x_{y_t}\to x$; they eventually avoid the finite exceptional set. Let
$t\downarrow0$ and use continuity. Only the top
coefficient remains, and the zero-mesh identity gives
$\sum_b\partial_b\operatorname{curl}(A_j)_{b\nu}=0$. The curl is also
closed by the Bianchi identity, hence is an on-shell Maxwell curvature.
The exact potential sequence represents this curvature as the curl of a
radial potential; uniqueness in radial gauge identifies that potential
with $A_j$. Thus $A_j$ lies in the image of the Maxwell potential map.
Continuity extends the assertion to $x=0$.

Now let a linear functional on the top Taylor layer annihilate that
solution-potential image. It vanishes on $A_j$ whenever the positive
force tests vanish on $A$. Thus its composition with projection onto the
top Taylor layer annihilates the common kernel of the force family.
Finite-dimensional duality puts this composed functional in the span of
the positive path forces. This is the force representation used to
construct the equation-factor read; no statement about the full
force-kernel in all Taylor layers is needed.

Write the equation-factor atom as one of grade $(j+2,1)$, and let the
spectator word have grade equal to the sum of the grades of its factors.
At the intermediate cutoff $j+2$, express the linearized equation factor
as a finite sum of transported link forces. Replacing each linearized
curvature by its actual slot read gives the forced factor with error in
slot filtration at least two. The spectator word is read by induction on
its length: for each factor, the safe-span lemma supplies slots avoiding
the force edge. The product rule then puts every error term at least one
slot above the spectator's own factor count. These reads have the desired
minimum slot count and preserve the one-link microscopic identity.

For the requested cutoff $d\geq j+2$, extend the intermediate Taylor
functional to the full jet and multiply by the unchanged spectator word.
The extension has two distinct errors. First, the discrepancy between
the constructed read and its own lifted Taylor polynomial has one more
slot factor than the target's second grade. Second, the lifted
linear functional agrees with the original equation atom on the
intermediate jet; its remaining Taylor tail has strictly higher first
grade while retaining the same slot count. Multiplying this tail by the
spectator word preserves the strict first-grade increase. Thus, writing
$w$ for the lifted polynomial and $v$ for the target generator, the
required comparison is the sum
\[
 w-v=\bigl(w-\operatorname{Taylor}_{\leq d}(\widetilde\ell\,a)\bigr)
       +\bigl(\operatorname{Taylor}_{\leq d}(\widetilde\ell\,a)-v\bigr).
\]
The first term is strictly later in the slot-count coordinate of the
grade stamp; the second is strictly later in the first-grade coordinate.
Both are therefore strictly above the target stamp. This two-part
comparison, together with the forced edge read, gives the required
positioned read for the generator. Finite sums and scalar multiples
extend it to the full homogeneous ideal element.

The reads in this lemma are raw reads; gauge averaging is done only after
the terms for the invariant label have been assembled.
\end{proof}

The conditional moment calculation in this argument is the one-link
Haar integration-by-parts identity. For a walk $w$ from the anchor to a
source edge $e$, the transported force is
\[
       \operatorname{Ad}_{\operatorname{travel}(U,0,w)}
                                  \operatorname{field}(U,e).
\]
Move the transport to the other slot using invariance of the color inner
product, then expand in a real orthonormal color basis $(E_a)$. Every
term has the form
\[
 \langle E_a,\operatorname{field}(U,e)\rangle\,S_a(U),
\]
where $S_a$ is the spectator times a coefficient depending only on the
walk transport. For the translated test at anchor $o$, choose the Micro
radius larger than the support radius of the walk and its source edge,
and choose the side threshold so these cover sites do not wrap. The test
spectator is then local outside that radius and is unchanged when only
the translated source link is varied. Haar integration by parts on that
link gives
\[
 \int \left(\partial_{E_a}S_a
       -\beta\langle E_a,\operatorname{field}(U,e)\rangle S_a\right)
            d\mu_\beta(U)=0.
\]
The derivative term vanishes; the derivative of the bare density is the
displayed force term. The envelope bound supplies an integrable majorant
for differentiating under the integral, so the identity also holds for
the bounded measurable spectators permitted in the Micro definition.
Summing over $a$ gives the force moment zero. This identity is valid for
either sign of $\beta$; the nonzero condition in Micro is used only when
dividing out the scalar factor in the moment normalization.

\paragraph{Scope of the factor lemma.}
The source's equation-ideal identity is the kernel statement used above.
The force calculation used is only the top-component implication proved
from the radial current limit: simultaneous vanishing of the positive
path-force tests puts the highest radial coefficient in the image of the
Maxwell potential map. The conditional one-link Haar identity is the
moment input defining the equation leg. The argument above proves the
monomial factorization, edge-avoiding spectator assembly, and two-part
filtration comparison from these facts.

\subsection{The all-grade campaign}

The previous lemmas give the exact pointwise producer needed for the
finite-dimensional assembly.

\begin{proposition}[Positioned equation-ideal read]\label{prop:marked-positioned}
For every pure-grade label $p$ in the Maxwell equation ideal, of grade
$m=(r,s)$ with $r\leq d$, there are a gauge-invariant function $f$ and a
polynomial $q$ such that $f$ has slot degree at least $s$, the action
derivatives through order $d$ of $f$ and $q$ agree, $q-p$ has grade
strictly above $m$, and $f$ is microscopic.
\end{proposition}
\begin{proof}
Use the exact ideal factorization to write $p$ as a finite sum of gap
atoms times spectator monomials. Apply the factor-read lemma to those
generators and add the resulting raw reads. Their polynomial targets sum
to $p$; their errors remain in the strictly higher-stamp subspace, and
the slot-degree and microscopic conditions pass to finite sums. The
resulting raw function need not yet be gauge invariant.

Now average the pair over the compact gauge group at the anchor. A
constant gauge transformation rotates the color components of each of
the finitely many fixed slot addresses among a finite basis. Thus the
orbit of this particular finite sum lies in a finite-dimensional span,
and its Haar integral is an ordinary finite-dimensional vector integral.
The label $p$ is gauge invariant, so the target of every rotated read is
still $p$; its Taylor correction remains in the same strictly higher
stamp subspace. Forms, slot filtration, and the microscopic moment
condition are linear and preserved by the integral. A based transported
slot under an arbitrary local gauge transformation changes only by the
adjoint action of that gauge transformation at its base point: the
interior transports cancel in pairs. Averaging over the anchor group
therefore makes the read invariant under the full local gauge action.
This produces the required $f$ and $q$.
\end{proof}

The proposition applies to the actual action, rather than to a formal substitute:
the slot germ is obtained from the link ODE above. Its one-factor
linearization is the oriented plaquette sum. The finite-generation formula
for the ideal is obtained degree by degree by expanding each polynomial
into monomials, applying the projection onto the Maxwell
solution-potential subspace, and splitting one
literal equation factor from the remaining slot coordinates. The
safe-span annihilator proof handles the edge avoidance needed to multiply
the spectators. These are the mathematical steps packaged by the
positioned read in the source.

\begin{theorem}[Campaign and imprint at fixed cutoff]
\label{thm:marked-campaign-producer}
Assume at cutoff $d$ there is a complex-linear map $S$ from labels to
link-field functions with unit value one, gauge-invariant values, and exact
action Taylor forms through order $d$. In addition, for each pure-grade
basis label of grade $(r,s)$, assume its value has a finite representation
as a linear combination of products of based curvature-slot color
coordinates, each product having at least $s$ factors. This is the powered
slot-degree condition used below. Then there is a complex-linear campaign
$F$ whose grade-$m$ value
has a read with slot degree at least the second grade, whose polynomial
correction has strictly larger stamp than the input label, and whose
equation-ideal values are microscopic. The campaign may be chosen to
commute with signed coordinate changes and the real structure.
\end{theorem}
\begin{proof}
We first construct the unsymmetrized campaign. The sound section itself
provides a linear map $X$ with
\[
 Xp=(p,S(p)),\qquad X1=(1,1).
\]
It gives the required action forms, gauge invariance, and slot degree for
each pure grade. It need not be microscopic on the equation ideal.

For a fixed grade $m$, restrict the input $p$ to $\operatorname{Grade}(m)$
and let
\[
 P_m=\operatorname{Grade}(m)\cap\mathcal I_d,
\]
where $\mathcal I_d$ is the finite-dimensional Maxwell equation ideal.
Choose a basis $(b_i)$ of $P_m$. For each $b_i$, apply
Proposition~\ref{prop:marked-positioned} and write its read as
$(q_i,f_i)$. Extend $b_i\mapsto(q_i,f_i)$ linearly on $P_m$. The result
is a payment map on the ideal part of this grade: forms, strict-stamp
correction, powered membership, and the Micro conclusion are all linear
conditions and hence pass through the basis expansion.

Choose a linear projection $L_m$ from the full label space onto $P_m$
which restricts to the identity on $P_m$. If $X_m$ is the map just
constructed on $P_m$ and $S_m$ is the sound-section map, set
\[
                 F_m(p)=X_m(L_mp)+S_m(p-L_mp).
\]
Both summands have the required grade-$m$ payment, so their sum does too.
If this grade-$m$ input $p$ lies in the equation ideal, grade preservation
of the ideal gives $L_mp\in P_m$ and $p-L_mp\in P_m$. Since $L_m$ is the
identity on $P_m$, applying $L_m$
to $p-L_mp$ gives zero and also $p-L_mp$; therefore $p-L_mp=0$.
Thus for ideal inputs the entire value is the microscopic marked value
$X_m(p)$. If $m=(0,0)$, the unit is not in the equation ideal: the ideal
has no nonzero constant component. Hence $L_m1=0$ and the sound map
supplies $F_m(1)=(1,1)$.

There are only finitely many relevant grades,
\[
             \{0,\ldots,d\}\times\{0,\ldots,d\}.
\]
Decompose $p=\sum_m p_m$ by its homogeneous grade and define
$F(p)=\sum_mF_m(p_m)$. On an input of pure grade $m$, every other grade
component vanishes, so the corresponding payment remains exact. On the
unit only grade $(0,0)$ contributes, so the start condition holds. This
is the claimed campaign.

Now symmetrize. Let $\mathcal S$ be the finite signed coordinate
permutation group, $\theta_s$ its action on labels, and $\phi_s$ its
joint action on a label/function pair. Define
\[
 F^{\mathrm{sgn}}=\frac1{|\mathcal S|}
       \sum_{s\in\mathcal S}\phi_s F\theta_{s^{-1}}.
\]
Signed changes preserve grade, the strictly higher-stamp subspace, the
slot-product filtration, the action-form identity, and the equation ideal.
They preserve the moment condition by changing variables in the finite
volume measure. Thus every summand remains a campaign and the average is
one. Reindexing by $s\mapsto ts$ shows
$F^{\mathrm{sgn}}\theta_t=\phi_tF^{\mathrm{sgn}}$.

Let $c$ be the joint real-conjugation involution on labels and
label/function pairs. It is conjugate-linear, commutes with the signed
action, and preserves each campaign condition. Set
\[
 F^{\mathrm{imp}}=\tfrac12\bigl(F^{\mathrm{sgn}}+
                           cF^{\mathrm{sgn}}c\bigr).
\]
The midpoint is complex-linear because the two conjugate-linear actions
on the source and target cancel. It remains a campaign; $c^2=1$ gives
real equivariance, while commutation of $c$ with the signed action retains
signed equivariance. Both averages fix the unit. This is the required
imprint.
\end{proof}

The exact algebraic assembly above uses no limit in volume and no
compactness argument. Its analytic input is the positioned read. The
source proves that input from the radial ODE germ, the finite Taylor
annihilator argument, the equation-ideal generation formula, and the
conditional-Haar microscopic identity described above. The section
input is supplied by the integrated section producer, as the following
corollary records.

\begin{corollary}[Actual fixed-cutoff imprint]\label{cor:marked-campaign-actual}
For every gauge group, holomorphic link action $H$, and cutoff $d$, the
source's actual section construction yields a campaign and an imprint at
that cutoff.
\end{corollary}
\begin{proof}
Use the local-sections construction of
Section~\ref{sec:local-sections}. For cutoff $d$ it gives the section
$T_dF$ and pose polynomials $P_\alpha$ with exact action forms through
order $d$, unit, and gauge covariance. Its pose certificate says that every
monomial in the pose polynomial of a grade-$(r,s)$ basis label has total
factor degree at least $s$. Expanding the finite pose polynomial writes each such basis value as
a finite sum of slot-coordinate products with at least $s$ factors. Thus
the powered slot-degree hypothesis of the theorem holds. Applying the
theorem gives a campaign; its finite signed and real averages give the
imprint.
\end{proof}
